\documentclass[12pt,a4paper]{article}
        \makeatletter
        \def\input@path{{../}}
        \makeatother
        \usepackage{almeria-fichas}

        \FichaMeta{Sesion 12}{De enunciado a tabla}{Bloque 3 · Datos, funciones y graficas}{Traducir situaciones verbales a tablas de datos claras y ordenadas.}{Procedimientos}{Comprender, Aplicar, Analizar}

        \begin{document}

        \FichaCabecera

        \begin{materialesbox}
        \begin{itemize}
    \item Ficha impresa
    \item Lapiz
    \item Regla
    \item Calculadora si se necesita
\end{itemize}
        \end{materialesbox}

        \begin{pasosbox}
        \begin{enumerate}
    \item Lee el enunciado dos veces.
    \item Rodea los datos que cambian y los que se repiten.
    \item Decide el nombre de cada columna antes de empezar.
    \item Ordena la tabla con cuidado.
\end{enumerate}
        \end{pasosbox}

        \begin{recordatoriobox}
        \begin{itemize}
    \item Toda tabla necesita encabezados claros.
    \item La variable independiente suele ir en la primera columna.
    \item Las unidades deben aparecer en el encabezado o junto al dato.
\end{itemize}
        \end{recordatoriobox}

        \begin{apoyobox}
        \begin{itemize}
    \item Puedes hacer primero un borrador muy sencillo de la tabla.
    \item Si un dato falta, deja hueco y relee el enunciado.
    \item No mezcles en la misma columna dos tipos de informacion.
\end{itemize}
        \end{apoyobox}

        \BloomSection{comprender}

\BloomExercise{comprender}{1}{Organizacion de datos}{Explica por que no conviene mezclar en una misma columna tiempo y temperatura.}{6}

\BloomExercise{comprender}{2}{Unidad visible}{Explica por que es mejor escribir \emph{Temperatura (ºC)} en lugar de solo \emph{Temperatura}.}{5}

\BloomSection{aplicar}

\BloomExercise{aplicar}{1}{Tabla desde texto}{Convierte en tabla este enunciado: \emph{Lunes 18 ºC, martes 20 ºC, miercoles 17 ºC, jueves 22 ºC}.}{7}

\BloomExercise{aplicar}{2}{Precio total}{Haz una tabla para el precio total de 1, 2, 3, 4 y 5 entradas si cada entrada cuesta 6 euros.}{7}

\BloomSection{analizar}

\BloomExercise{analizar}{1}{Partes y relaciones}{Analiza un ejemplo relacionado con \emph{De enunciado a tabla}. Separa sus partes, indica cuales son relevantes y explica como se relacionan entre si.}{8}

\BloomExercise{analizar}{2}{Detecta errores o diferencias}{Compara dos respuestas, dos representaciones o dos procedimientos relacionados con esta sesion. Analiza sus diferencias y explica que errores o mejoras detectas.}{8}

        \begin{evidenciabox}
        \begin{itemize}
    \item Pasar de texto a tabla sin perder informacion.
    \item Nombrar bien las columnas y las unidades.
    \item Ordenar datos de forma legible.
\end{itemize}
        \end{evidenciabox}

        \end{document}
